\subsection{The singular support gate requires three distinct variables}

\begin{proposition}[P-F27: extremal determinant order for graph supports]
\label{prop:cubicgraphsupport}
Let $n\ge2$, and let $F\in\mathbb C[x_1,\ldots,x_n]$ be a
cubic whose nonzero monomials all have the forms $x_i^3$ or
$x_i^2x_j$ with $i\ne j$. Let $B$ be the integer matrix whose
columns are their exponent vectors. Assume that the vector
$\mathbf1=(1,\ldots,1)^t$ belongs to the rational column span
of $B$. This hypothesis holds, in particular, if $F$ is
$SL_n$-semistable.

Suppose a diagonal matrix $A$ of finite two-power order fixes
$F$ exactly. Then the order of $\det A$ is at most $2^{n-1}$.
If that order is $2^{n-1}$, the support is, after a permutation
of variables, exactly
\[
 x_1^2x_2,\ x_2^2x_3,\ldots,\ x_{n-1}^2x_n,\ x_n^3.
 \tag{P-F27a}
\]
All their coefficients are nonzero, the matrix $A$ itself has
order $2^{n-1}$, and a diagonal change of variables turns $F$
into the smooth chain cubic $F_n$.

Consequently, for $n\ge4$, every nonzero fixed quotient point
in the discriminant boundary of P-F26 has the following
property: in any eigenbasis for any finite two-power
exact-form lift with determinant of order $2^{n-1}$, its
closed-orbit representative contains a nonzero monomial
$x_ix_jx_k$ with $i,j,k$ pairwise distinct.
\end{proposition}

\begin{proof}
We first explain the rational-span assumption. If a
positive-degree $SL_n$-invariant polynomial in the
coefficients of $F$ does not vanish at $F$, at least one of
its coefficient monomials has nonzero value. Invariance
under the diagonal determinant-one torus means that the
sum of the cubic exponent vectors, counted with that
coefficient monomial's nonnegative integer multiplicities,
has all coordinates equal. The common coordinate is
positive, because the coefficient degree is positive.
Dividing by it expresses $\mathbf1$ as a nonnegative
rational combination of columns of $B$. Only membership
in their rational span will be used below.

Write the eigenvalues of $A$ as $\alpha_1,\ldots,\alpha_n$.
They lie in one finite cyclic two-group. The map
\[
 \psi:\mathbb Z^n\longrightarrow\mathbb C^*,\qquad
 (u_i)\longmapsto\prod_i\alpha_i^{u_i}
 \tag{P-F27b}
\]
annihilates the lattice $L$ generated by the columns of $B$,
since every supported monomial is fixed. Thus it factors
through $\mathbb Z^n/L$, and
$\psi(\mathbf1)=\det A$. This is the precise lattice
interpretation of the determinant character; the order
of the determinant need not equal the order of $A$.

Make a graph on the variables. A supported monomial
$x_i^2x_j$ gives an edge with coefficient two at vertex
$i$ and coefficient one at vertex $j$; a cube gives a loop
with coefficient three. Distinct oriented monomials may
give two edges on the same unordered pair. Decompose the
graph into connected components, ignoring loops for
connectivity. The exponent lattice and rational span
split over these components. In a component with $v$
vertices, therefore, its all-ones vector $\mathbf1_v$
belongs to the rational span of its support columns.
The determinant is the product of the images of these
component vectors under $\psi$.

We prove that each component contributes a determinant
character of order at most $2^{v-1}$, with a stronger
bound $2^{v-2}$ in the rank-deficient case when $v\ge2$.
An isolated vertex must have a cube, since otherwise
$\mathbf1_v$ would not belong to the column span. Its
relation is $3e_i=0$, and every character into a two-group
is trivial on $e_i$. This handles $v=1$.

For $v\ge2$, choose a spanning tree of the underlying
connected graph, and let $T$ be the $v\times(v-1)$ matrix
of its support columns. Its columns are independent:
removing a leaf removes a row with exactly one nonzero
entry, equal to one or two, and induction proves rank
$v-1$. Hence the full component matrix has rank either
$v-1$ or $v$.

\emph{The rank-deficient case.}
Suppose its rank is $v-1$. Then its rational column span
is that of $T$, so $Tc=\mathbf1_v$ has a unique rational
solution $c$. For every vertex $r$, let $T^{(r)}$ be the
square matrix obtained by deleting row $r$. Repeatedly
remove leaves other than $r$. This proves
\[
 D_r:=\det T^{(r)}=\pm2^{e_r},\qquad 0\le e_r\le v-1.
 \tag{P-F27c}
\]
Cramer's rule applied after deleting row $r$ shows that
the least positive integer $q$ with $qc$ integral divides
$|D_r|$. It therefore divides the greatest common divisor
of all these maximal minors. Also
$q\mathbf1_v=T(qc)$ lies in the tree lattice. Equivalently,
the torsion class of $\mathbf1_v$ in the tree cokernel has
order $q$, and its image in the full support cokernel has
order dividing $q$. This proves the needed maximal-minor
assertion directly, without assuming full rank.

The signed cofactor vector of $T$ is a nonzero vector
in the kernel of $T^t$. Along an edge with coefficients
two and one, its two endpoint coordinates consequently
satisfy $2u_i+u_j=0$. Their absolute values differ by a
factor of two. Since their absolute values are the
$|D_r|$, two adjacent exponents in (P-F27c) differ by one.
As all these exponents are at most $v-1$, their minimum
is at most $v-2$. Thus
\[
 q\mid 2^{\min_r e_r},\qquad \min_r e_r\le v-2.
 \tag{P-F27d}
\]
The component determinant character has order at most
$2^{v-2}$, as asserted. Notice that free directions in
the cokernel may allow the matrix eigenvalues themselves
to have much larger order; the span hypothesis makes
$\mathbf1_v$ a torsion class and is what controls its
character.

\emph{The full-rank case.}
Suppose instead that the component matrix has rank $v$.
One support column outside the span of $T$ extends its
columns to a square nonsingular matrix $C$. Its support
graph is connected with $v$ edges and $v$ vertices, and
has exactly one cycle, possibly a loop or a double edge.
The finite cokernel of $C$ has order $|\det C|$; one may
also use the adjugate identity to see directly that
$\det C$ kills every cokernel class. Since its image under
$\psi$ is a two-group, the component determinant order
is at most $2^{v_2(\det C)}$.

Let the unique cycle have length $l$. Expanding the
matrix successively along leaf rows contributes one
factor of either one or two for each of its $v-l$
tree edges. If $l=1$, the remaining loop contributes
three, and hence
\[
 v_2(\det C)\le v-1.
 \tag{P-F27e}
\]
If $l\ge2$, the cycle determinant has exactly two
permutation terms. For some $0\le a\le l$ they give,
up to an overall sign,
\[
 2^a+(-1)^{l-1}2^{l-a}.
 \tag{P-F27f}
\]
If $a=l-a$, then $l$ is even and the determinant is zero,
which is excluded. Otherwise its two-adic valuation is
$\min(a,l-a)\le\lfloor(l-1)/2\rfloor$. Therefore
\[
 v_2(\det C)
 \le v-l+\left\lfloor\frac{l-1}{2}\right\rfloor
 \le v-2\qquad(l\ge2).
 \tag{P-F27g}
\]
This proves the component bound in every case, including
multiple edges and rank deficiency.

Because all component characters have two-power order,
their product has order at most their largest order.
Every component has at most $n$ vertices. This proves
$\operatorname{ord}(\det A)\le2^{n-1}$.

Now suppose equality holds. At least one component must
have $n$ vertices, since every smaller component has
bound at most $2^{n-2}$. Thus the support graph is
connected. It has full rank by (P-F27d), and the chosen
unicyclic basis must have a loop by (P-F27g). Moreover,
equality in (P-F27e) requires every leaf expansion factor
to be two. Root its remaining tree at the loop vertex.
Each basis edge therefore has coefficient two at its
child and one at its parent. The eigenvalue relations are
\[
 \alpha_{\mathrm{root}}^3=1,\qquad
 \alpha_{\mathrm{parent}}=\alpha_{\mathrm{child}}^{-2}.
 \tag{P-F27h}
\]
Since all eigenvalues have two-power order, the root
eigenvalue is one. A vertex at depth $d$ then has
eigenvalue order at most $2^d$. If the maximum tree
depth is $h$, the order of every eigenvalue, and therefore
of their product, is at most $2^h$. Equality for the
determinant forces $h\ge n-1$. A rooted tree with $n$
vertices has $h\le n-1$, with equality only for a path
rooted at an endpoint. Thus this basis is the chain
support in (P-F27a).

The deepest eigenvalue must have order $2^{n-1}$: all
other vertices have smaller depth, and a product of
elements of smaller two-power order could not have the
required determinant order. Write this eigenvalue as
$\xi$, and order the path from its deepest vertex to the
root. Then the entire eigenvalue list is
\[
 \xi,\ \xi^{-2},\ \xi^4,\ldots,
 \xi^{(-2)^{n-1}}=1,
 \qquad \operatorname{ord}(\xi)=2^{n-1}.
 \tag{P-F27i}
\]
In particular, $A$ has this same order, although this
was not assumed.

These eigenvalues permit no additional cubic support.
Indeed their nonzero residue weights modulo $2^{n-1}$
have distinct two-adic valuations. In a zero sum of
three weights, the least nonzero valuation cannot occur
once, and cannot occur three times because three is
odd. It occurs twice at the same vertex; the third
weight must be its negative double, the next chain
weight. The all-zero case gives the root cube. Thus the
fixed cubic monomials are exactly (P-F27a). They are all
present, since they were the columns of the chosen basis.
Their nonzero coefficients can be made one by first
rescaling the root variable using a cube root and then
successively rescaling its predecessors using square
roots. Hence $F$ is the smooth chain cubic up to a
diagonal substitution, by the smoothness proof in P-F19.

Finally, P-F26 supplies a finite two-power lift with
primitive determinant for every nonzero fixed quotient
point. Its closed-orbit representative is semistable,
so it satisfies the span hypothesis in every coordinate
basis. If that representative is singular and any such
lift eigenbasis had only graph-support monomials, the
result just proved would identify it with a smooth
chain. This contradiction proves the last assertion.
\end{proof}

\paragraph{Connected graph part, even with additional triple monomials.}
The same equality classification holds for an arbitrary cubic satisfying
the rational-span hypothesis if its graph made from its square-linear
monomials is connected; cubes are retained as loops. Thus the presence
of a three-distinct-variable monomial by itself is not enough to produce
a competitor. In the discriminant boundary of P-F26, this graph must be
disconnected in every finite-lift eigenbasis.

Here is the additional proof. Let $G$ be the matrix consisting only
of the cube and square-linear support columns. Since its graph is
connected, its rank is $n$ or $n-1$. If it has rank $n$, its graph
subform satisfies the span hypothesis by itself. The proposition
therefore gives the chain eigenvalue list (P-F27i), which permits no
additional triple monomials. The full form is the chain.

Suppose instead that $G$ has rank $n-1$, and take the spanning-tree
matrix $T$ used above. Its signed cofactor vector has the form
\[
 u_i=\varepsilon 2^e(-2)^{h_i},\qquad
 \varepsilon\in\{1,-1\},\quad e\ge0,\quad
 h_i\ge0,\quad \min_i h_i=0.
 \tag{P-F27j}
\]
Indeed its coordinates are signed powers of two and the relation
$2u_i+u_j=0$ along each tree edge fixes both the change in exponent
and the change in sign. Connectivity makes the sign formula in
(P-F27j) consistent. Put $H=\max_i h_i$. The tree-minor calculation
gives $e+H\le n-1$.

If every triple support column lies in the rational span of $T$,
then so does $\mathbf1$, and the rank-deficient argument already
bounds the determinant order by $2^{n-2}$. Otherwise choose a triple
column $\beta=e_i+e_j+e_k$ outside that span, with $i,j,k$ distinct.
Cofactor expansion gives a nonsingular square matrix with
\[
 \det(T,\beta)=\pm 2^e
       \bigl((-2)^{h_i}+(-2)^{h_j}+(-2)^{h_k}\bigr).
 \tag{P-F27k}
\]
A nonzero sum of three powers of $-2$ has two-adic valuation strictly
less than the largest exponent occurring, unless all three exponents
are equal. To verify this, a unique smallest exponent gives exactly
that valuation. If two smallest exponents are $r$ and the third is
$t>r$, the sum is $(-2)^r(2+(-2)^{t-r})$: it vanishes when $t=r+1$,
and has valuation $r+1<t$ when $t\ge r+2$. Three equal exponents
$r$ give valuation $r$. Consequently the valuation in (P-F27k) is
at most $e+H$, with equality possible only if all three chosen
vertices have height $H$.

If $e+H\le n-2$, this already contradicts determinant order
$2^{n-1}$, since the determinant of the square support submatrix
annihilates the determinant character. In the only remaining case,
$e+H=n-1$, choose a vertex $r$ of height $H$. Its deleted-row tree
minor has absolute value $2^{n-1}$. Root the tree at $r$ and perform
the leaf expansion for this minor. Every one of its $n-1$ factors
must be two. Thus each tree edge has coefficient two at its child
and one at its parent, and the cofactor relation says
$h_{\rm parent}=h_{\rm child}+1$. The root is the unique vertex of
height $H$. Three distinct vertices cannot all have that height.
Therefore the valuation in (P-F27k) is again at most $n-2$, the
same contradiction. This completes the connected-graph extension.

\paragraph{Scope of the reduction.}
The proof treats graph supports and, by the extension just proved,
arbitrary supports with connected graph part. It does not classify
all singular semistable cubics. It does not infer
that the determinant order bounds an arbitrary finite lift
before imposing the extremal graph-support case. Three
distinct-variable monomials are exactly the remaining
cubic support type, and the unresolved case has a disconnected
graph part. Thus P-F27 removes the rank-deficiency and
larger-order-lift obstructions for the stated classes, while
leaving the disconnected hypergraph support problem unresolved by
this proposition alone. P-F30 subsequently resolves that remaining
case by component contraction and a determinant bound, and P-F31
uses it to control the global fixed support and ordinary pole.
No historical priority is claimed for this elementary
lattice and graph reduction.
